\begin{lemma}[the pair-of-pants defect class]
\hypertarget{tag-03IC}{}\label{wbII:lem:picard-crossed-defect}
\StatementRecord{wbII:lem:picard-crossed-defect}{proved}{03IC}{\SeriesRef{wbII:eq:W:5.17j}; \cite[Chapter~III, \S1, Exercise~2, p.~60; Chapter~IV, \S2, pp.~88--89]{Brown1982}}

\ Choose the compact pair-of-pants core $P\subset B^\circ$ and
the zero-section whiskers used in \CompendiumRef{I}{mapsI:top:prop:geometric-pi1}, and put
$J_P=f^{-1}(P)$.  Orient its boundary meridians clockwise.  Write
$\rho_0$ for the projection of the zero-section cusp meridian
$\widetilde a_0$.  The standard family uses the tautological cusp
gluing, whose affine regluing class is $\mu=0$ (equivalently
$\ell_0=\gamma(\mu)=0$); hence the flat-marked cusp lift has no vertical
translation.  The base loops and their chosen lifts therefore satisfy
\[
 \rho_1\rho_2\rho_0=1,
 \qquad
 \widehat\rho_1\widehat\rho_2\widetilde a_0=1.
\tag{W.23}\hypertarget{tag-03GQ}{}\label{wbII:eq:W:5.17g}
\]
For $L\in\operatorname{Pic}(X)$, write
$E=c_1(L|_F)\in\bigwedge^2V$.  Then $E\in(\bigwedge^2V)^G$, and there is a
crossed homomorphism
\[
 z:\pi_1(P)\longrightarrow V_{\mathbb R},\qquad
 z(gh)=z(g)+T_gz(h),
\tag{W.24}\hypertarget{tag-03GR}{}\label{wbII:eq:W:5.17h}
\]
whose class records
$c_1(L|_{J_P})-[\Omega_E]$.  Here $\Omega_E$ is the closed, translation-
invariant vertical two-form induced by $E$ in the flat marking.  If
$d_i=z(\rho_i)$ and a bar denotes the image in
$V_G\otimes\mathbb R$, then
\[
 \bar d_1+\bar d_2+\bar d_0=0.
\tag{W.25}\hypertarget{tag-03GS}{}\label{wbII:eq:W:5.17i}
\]
The three barred classes, and the number
$\langle d_0,\widehat\delta\rangle$, do not depend on the chosen crossed-
homomorphism representative.  More precisely, use the base-first convention
for the edge injection in \eqref{wbII:eq:W:5.17j}.  For any oriented loop $\rho\subset P$
and any $A_\rho$-invariant $\nu\in\Lambda$, its flat-marked sweep satisfies
\[
 \bigl\langle z(\rho),\nu\bigr\rangle
 =-\bigl\langle c_1(L|_{J_P})-[\Omega_E],
 a_*([S^1_\nu]\times[\widetilde\rho])\bigr\rangle.
\tag{W.26}\hypertarget{tag-03GU}{}\label{wbII:eq:W:5.17j'}
\]
Here $a$ is the fibrewise translation action and $\widetilde\rho$ is the
flat-marked lift of $\rho$.
\end{lemma}



\begin{proof}

\begin{enumerate}[label={$\langle1\rangle$\arabic*.},leftmargin=*]
\item The flat marking identifies the pullback of $J_P$ to the universal cover of
$P$ with
$\widetilde P\times(\Lambda_{\mathbb R}/\Lambda)$.  Since $E$ is the
restriction of a class on the total space, it is fixed by every monodromy
$T_g$.  The constant vertical form representing $E$ is therefore invariant
under the deck action and descends to the asserted form $\Omega_E$ on
$J_P$.

\item The pair of pants retracts onto a graph.  Hence the Serre spectral sequence
of the marked real torus bundle has only columns zero and one and gives the
short exact edge sequence
\[
0\longrightarrow H^1(P;V_{\mathbb R})
 \longrightarrow H^2(J_P;\mathbb R)
 \longrightarrow H^0(P;\textstyle\bigwedge^2V_{\mathbb R})
 \longrightarrow0.
\tag{W.27}\hypertarget{tag-03GT}{}\label{wbII:eq:W:5.17j}
\]
Both $c_1(L|_{J_P})$ and $[\Omega_E]$ map to $E$, so their difference is
the image of a unique class in $H^1(P;V_{\mathbb R})$.  The standard group-
cohomology model for a graph identifies this group with crossed
homomorphisms modulo principal crossed homomorphisms
\cite[Chapter~III, \S1, Exercise~2, p.~60; Chapter~IV, \S2, pp.~88--89]{Brown1982}; this gives \eqref{wbII:eq:W:5.17h}.

\item For completeness, the slant convention in \eqref{wbII:eq:W:5.17j'} is part of this
identification, not an additional geometric assertion.  Give an oriented
edge of a graph spine of $P$ a coordinate $s$ increasing with its based
loop and represent the corresponding $V_{\mathbb R}$-valued cellular
one-cochain by a base-first form $ds\wedge\beta_s$.  On the sweep of an
$A_\rho$-invariant fibre circle $S^1_\nu$, Fubini's formula gives
\[
\int_{S^1_s\times S^1_\nu}ds\wedge\beta_s
 =\int_{S^1_s}\langle\beta_s,\nu\rangle,ds
 =\langle z(\rho),\nu\rangle.
\]
The sweep in \eqref{wbII:eq:W:5.17j'} is written in the fibre-first order
$[S^1_\nu]\times[\widetilde\rho]$, which differs from this base-first
cross-product orientation by $(-1)^{1\cdot1}=-1$; this is the displayed
minus sign.
Changing the cellular representative by a principal cochain changes both
sides by the same boundary term, which vanishes because $A_\rho\nu=\nu$.
This proves \eqref{wbII:eq:W:5.17j'}.

\item Applying $z$ to \eqref{wbII:eq:W:5.17g} gives
\[
 d_1+T_1d_2+T_1T_2d_0=0.
\]
All $T_g$ act trivially on the global coinvariants, proving \eqref{wbII:eq:W:5.17i}.  A
change of representative has the form
$z(g)\mapsto z(g)+(T_g-I)a$.  Its image in $V_G\otimes\mathbb R$ is zero.
Moreover the displayed matrices in
\CompendiumRef{I}{wb:prop:monodromy-certificate} fix $\widehat\delta$ under every dual
action $A_g$.  Therefore
\[
 \langle(T_g-I)a,\widehat\delta\rangle
 =\langle a,(A_g^{-1}-I)\widehat\delta\rangle=0,
\]
which proves the final assertion.

\end{enumerate}
\end{proof}